%% kapa: exam=1 round=37 year=2026 period=2 form=A subject=accounting subject_ko=회계학 seq=1 number=41
%% source: Q-Net 감정평가사 기출문제 (cst00309) artlSeq=5258924 file=raw/1008/5258924/2026년 제37회 감정평가사 1차 2교시.pdf
%% conversion: kordoc 4.24.51 -> tools/gen1.py (auto); review=proofread-sonnet; flags=none
\begin{question}{41}
재무보고를 위한 개념체계에 관한 설명으로 옳지 않은 것은?
\begin{choices}
  \item 개념체계는 특정 거래나 다른 사건에 적용할 회계기준이 없거나 회계기준에서 회계정책 선택이 허용되는 경우에 재무제표 작성자가 일관된 회계정책을 개발하는 데 도움을 준다.
  \item 현재 및 잠재적 투자자, 대여자와 그 밖의 채권자는 일반목적재무보고서의 대상이 되는 주요 이용자이다.
  \item 목적적합한 재무정보는 이용자들의 의사결정에 차이가 나도록 할 수 있으며, 예측가치와 확인가치 또는 둘 모두가 있으면 그 재무정보는 의사결정에 차이가 나도록 할 수 있다.
  \item 인식은 자산, 부채, 자본, 수익 또는 비용과 같은 재무제표 요소 중 하나의 정의를 충족하는 항목을 재무상태표나 재무성과표에 포함하기 위하여 포착하는 과정이다.
  \item 보고기업은 반드시 법적 실체이어야 하고 단일의 실체이거나 어떤 실체의 일부일 수 있으며, 둘 이상의 실체로 구성될 수 있다.
\end{choices}
\end{question}
\answer{5}
